\documentclass[10pt]{article}
\usepackage[margin=1in]{geometry}
\usepackage[T1]{fontenc}
\usepackage{lmodern,microtype,amsmath,amssymb,amsthm}
\theoremstyle{definition}
\newtheorem{assumption}{Assumption}
\begin{document}
\section*{Derivative Convolutional Perceptron Separation}
We seek to separate samples $b_0,\ldots,b_n\in\mathbb{R}^d$, where $d\geq1$ and $n\geq0$ are integers, using parameters $\theta=((a_1,a_2),w)\in\mathbb{R}^2\times\mathbb{R}^d$. With $S$ denoting the cyclic coordinate shift, the margin is $m_i(\theta)=a_1\langle b_i,w\rangle+a_2\langle Sb_i,w\rangle$. Each sample comprises a bounded, nonzero constant background and a small centered defect; the normalized cyclic differences of the defects share a unit separator with positive margin. The defect scale satisfies a quantitative smallness condition relative to the background bounds and derivative margin. The oracle returns $m_i(\theta)$ and its parameter derivative at a queried pair $(i,\theta)$. The goal is an output $\bar{\theta}$ with $m_i(\bar{\theta})>0$ for every sample and every realization of algorithm randomness.

Equip the feature space $\mathbb{R}^d$ with its Euclidean inner product and norm. Index coordinates by $j=0,\ldots,d-1$, and write
\[
(Sv)_j=v_{(j-1)\bmod d},
\qquad
u=\frac{1}{\sqrt d}(1,\ldots,1).
\]
The task is strict feasibility of all sample margins in the parameter space $\mathbb{R}^2\times\mathbb{R}^d$. The algorithm receives $d,n,g,R,\mu,\kappa$, where
\[
0<g\leq R,\qquad \mu>0,\qquad 0<\kappa\leq1;
\]
the samples and their geometric witnesses are accessible only through the oracle.

Our assumptions:
\begin{assumption}[Background and derivative separation]
There exist scalars $c_0,\ldots,c_n\in\mathbb{R}$, defects $v_0,\ldots,v_n\in\mathbb{R}^d$, and a common separator $s\in\mathbb{R}^d$ such that $\|s\|=1$ and, for every $i=0,\ldots,n$,
\[
\begin{aligned}
b_i&=\frac{c_i}{\sqrt2}u+\mu v_i,\
g&\leq |c_i|\leq R,\
\langle u,v_i\rangle&=0,\qquad \|v_i\|\leq\frac12,\
\left\langle\frac{v_i-Sv_i}{\sqrt2},s\right\rangle&\geq\kappa.
\end{aligned}
\]
\end{assumption}

\begin{assumption}[Contrast scale]
The contrast and background scales satisfy
\[
\frac{\mu}{R}\leq
\frac{\kappa^{20}}{2^{400}}\left(\frac{g}{R}\right)^{24}.
\]
\end{assumption}

At any query $(i,\theta)\in\{0,\ldots,n\}\times(\mathbb{R}^2\times\mathbb{R}^d)$, the oracle returns the exact scalar margin and derivative direction
\[
\begin{aligned}
m_i(\theta)&=a_1\langle b_i,w\rangle
                 +a_2\langle Sb_i,w\rangle,\
h_i(\theta)&=\bigl((\langle b_i,w\rangle,\langle Sb_i,w\rangle),
                    a_1b_i+a_2Sb_i\bigr).
\end{aligned}
\]
Here $h_i(\theta)$ represents the parameter derivative under the Euclidean coordinate pairing. Initialization is chosen by the algorithm. After the allotted oracle calls, its output $\bar{\theta}\in\mathbb{R}^2\times\mathbb{R}^d$ must satisfy
\[
m_i(\bar{\theta})>0
\qquad\text{for every }i=0,\ldots,n,
\]
for every realization of its randomness.
\end{document}
